% Part:sets-functions-relations
% Chapter: size-of-sets
% Section: non-enumerability

\documentclass[../../../include/open-logic-section]{subfiles}

\begin{document}

\olfileid{sfr}{siz}{nen}

\olsection{\printtoken{S}{nonenumerable} సమితులు}

\begin{editorial}
  ఈ విభాగం \olref[enm]{sec}లోని నిర్వచనాన్ని వాడి $\Bin^\omega$,
  $\Pow{\PosInt}$ లెక్కించలేనివని నిరూపిస్తుంది. ఇది
  \olref[enm-alt]{sec}లోని రూపం కంటే కొంచెం ప్రాథమికంగా, మరింత
  వివరంగా ఉండేలా రూపొందించబడింది.
\end{editorial}

$\PosInt$ ధన పూర్ణ సంఖ్యల సమితి వంటి కొన్ని సమితులు అనంతమైనవి.
ఇప్పటివరకు చూసిన అనంత సమితుల ఉదాహరణలన్నీ \tetoken{లెక్కించదగిన}{!!{enumerable}}. అయితే ఈ
లక్షణం లేని అనంత సమితులూ ఉన్నాయి. అటువంటి సమితులను
\emph{\tetoken{లెక్కించలేని}{!!{nonenumerable}}} అంటారు.

మొదట, \tetoken{లెక్కించలేని}{!!{nonenumerable}} సమితులు ఉండటమే ఆశ్చర్యంగా అనిపించవచ్చు.
ఏ \tetoken{లెక్కించదగిన}{!!{enumerable}} సమితి $A$కైనా \tetoken{ఒక సంగ్రస్త}{!!a{surjective}} ప్రమేయం
$f \colon \PosInt \to A$ ఉంటుంది. ఒక సమితి \tetoken{లెక్కించలేని}{!!{nonenumerable}} అయితే
అలాంటి ప్రమేయం ఉండదు. అంటే, $\PosInt$లోని అనంత సంఖ్యలో
\tetoken{మూలకాలను}{!!{element}s} $A$కు పంపే ఏ ప్రమేయమూ $A$ మొత్తాన్ని నింపలేదు.
కాబట్టి $A$లో, అనంత సంఖ్యలో ఉన్న ధన పూర్ణ సంఖ్యలకంటే ``ఎక్కువ''
\tetoken{మూలకాలు}{!!{element}s} ఉన్నాయి.

ఒక సమితి \tetoken{లెక్కించలేని}{!!{nonenumerable}} అని ఎలా నిరూపించాలి? అలాంటి సంగ్రస్త
ప్రమేయం ఏదీ ఉండజాలదని చూపాలి. సమానార్థకంగా, $A$లోని మూలకాలను ఒక
దిశలో అనంతంగా సాగే జాబితాలో లెక్కించలేమని చూపాలి. ఇందుకు అత్యుత్తమ
పద్ధతి, $A$లోని \tetoken{మూలకాలతో}{!!{element}s} చేసిన ప్రతి జాబితా కనీసం ఒక మూలకాన్ని
వదిలివేస్తుందని, అంటే ఏ ప్రమేయం $f\colon \PosInt \to A$ కూడా
\tetoken{సంగ్రస్త}{!!{surjective}} కాదని చూపడం. కాంటర్ \emph{వికర్ణ పద్ధతి}తో దీన్ని
చేయవచ్చు. $A$లోని \tetoken{మూలకాలు}{!!{element}s} జాబితా $x_1$, $x_2$, \dots ఇచ్చినప్పుడు,
నిర్మాణం వల్లనే ఆ జాబితాలో ఉండజాలని $A$ యొక్క మరో మూలకాన్ని నిర్మిస్తాం.

మన మొదటి ఉదాహరణ ఖాళీలు లేని అనంత $0$, $1$ క్రమాలన్నిటి సమితి
$\Bin^\omega$.

\begin{thm}
\ollabel{thm:nonenum-bin-omega}
$\Bin^\omega$~\tetoken{లెక్కించలేని}{!!{nonenumerable}}.
\end{thm}

\begin{proof}
విరోధం కోసం $\Bin^\omega$ \tetoken{లెక్కించదగిన}{!!{enumerable}} అని అనుకుందాం; అంటే,
$\Bin^\omega$లోని \tetoken{మూలకాలు}{!!{element}s} అన్నిటి జాబితా $s_{1}$, $s_{2}$,
$s_{3}$, $s_{4}$, \dots{} ఉందనుకుందాం. ఈ ప్రతి $s_i$ స్వయంగా
$0$, $1$ల అనంత క్రమం. ఈ జాబితాలోని $i$వ క్రమంలోని $j$వ మూలకాన్ని
$s_i(j)$ అందాం. అప్పుడు $i$వ క్రమం $s_i$ ఇలా ఉంటుంది:
\[
s_i(1), s_i(2), s_i(3), \dots
\]

ఈ జాబితానూ, అందులోని ప్రతి క్రమం $s_i$ యొక్క మూలకాలనూ ఒక పట్టికలో
అమర్చవచ్చు:
\[
\begin{array}{c|c|c|c|c|c}
& 1 & 2 & 3 & 4 & \dots \\\hline
1 & \mathbf{s_{1}(1)} & s_{1}(2) & s_{1}(3) & s_1(4) & \dots \\\hline
2 & s_{2}(1)& \mathbf{s_{2}(2)} & s_2(3) & s_2(4) & \dots \\\hline
3 & s_{3}(1)& s_{3}(2) & \mathbf{s_3(3)} & s_3(4) & \dots \\\hline
4 & s_{4}(1)& s_{4}(2) & s_4(3) & \mathbf{s_4(4)} & \dots \\\hline
\vdots & \vdots & \vdots & \vdots & \vdots & \mathbf{\ddots}
\end{array}
\]
పక్కన ఉన్న సంఖ్యలు $s_1$, $s_2$, \dots జాబితాలో క్రమం సంఖ్యను
ఇస్తాయి; పైనున్న సంఖ్యలు ఒక్కొక్క క్రమంలోని \tetoken{మూలకాలకు}{!!{element}s} పేర్లు.
ఉదాహరణకు $s_{1}(1)$ అనేది $s_{1}$ క్రమంలోని మొదటి \tetoken{మూలకం}{!!{element}} అయిన
$0$ లేదా $1$కు పేరు; తదితరంగా.

ఇప్పుడు ఈ జాబితాలో ఉండే అవకాశమే లేని $0$, $1$ల అనంత క్రమం
$\overline{s}$ను నిర్మిస్తాం. $\overline{s}$ నిర్వచనం $s_1$, $s_2$,
\dots జాబితాపై ఆధారపడుతుంది. $0$, $1$ల అనంత క్రమాల ఏ అనంత జాబితా
ఇచ్చినా, దానిలో కనిపించదని హామీ ఉన్న అనంత క్రమం $\overline{s}$ను
అది ఉత్పత్తి చేస్తుంది.

$\overline{s}$ను నిర్వచించడానికి దాని \tetoken{మూలకాలు}{!!{element}s} అన్నిటినీ, అంటే
ప్రతి $n \in \PosInt$కు $\overline{s}(n)$ను పేర్కొంటాం. పై పట్టికలో
వికర్ణం వెంబడి చదివి (అందుకే ``వికర్ణ పద్ధతి'' అనే పేరు), ప్రతి
$1$ను $0$గానూ, ప్రతి $0$ను $1$గానూ మార్చి దీన్ని చేస్తాం. మరింత
అమూర్తంగా, వికర్ణంలోని $n$వ \tetoken{మూలకం}{!!{element}} $s_n(n)$ విలువ $1$ లేదా $0$
అవడాన్ని బట్టి $\overline{s}(n)$ను $0$ లేదా $1$గా నిర్వచిస్తాం:
\[
\overline{s}(n) =
\begin{cases}
1 & \text{$s_{n}(n) = 0$ అయితే}\\
0 & \text{$s_{n}(n) = 1$ అయితే}.
\end{cases}
\]
సందర్భాలవారీ నిర్వచనం కంటే సూత్రం నచ్చితే,
$\overline{s}(n) = 1 - s_n(n)$గా కూడా నిర్వచించవచ్చు.

వికర్ణంలో కనిపించే $0$, $1$ల క్రమానికి ఇది ప్రతిబింబ క్రమమే కనుక,
$\overline{s}$ స్పష్టంగా $0$, $1$ల అనంత క్రమం. అందుచేత
$\overline{s}$ అనేది $\Bin^\omega$ యొక్క \tetoken{ఒక మూలకం}{!!a{element}}. కానీ అది
$s_1$, $s_2$, \dots జాబితాలో ఉండజాలదు. ఎందుకు?

అది జాబితాలోని మొదటి క్రమం $s_1$ కాలేదు, ఎందుకంటే మొదటి
\tetoken{మూలకంలో}{!!{element}} $s_1$కు భిన్నంగా ఉంటుంది. $s_1(1)$ ఏదైనా,
$\overline{s}(1)$ను దానికి వ్యతిరేకంగా నిర్వచించాం. అది రెండవ క్రమం
కాలేదు, ఎందుకంటే రెండవ మూలకంలో $\overline{s}$, $s_2$కు భిన్నంగా
ఉంటుంది: $s_2(2)$ విలువ $0$ అయితే $\overline{s}(2)$ విలువ $1$;
అదే విధంగా తిరుగునూ. ఇలాగే కొనసాగుతుంది.

మరింత కచ్చితంగా: $\overline{s}$ జాబితాలో ఉంటే, ఏదో ఒక $k$కు
$\overline{s} = s_{k}$ అవుతుంది. రెండు క్రమాలు ప్రతి స్థానంలో
ఏకీభవిస్తేనే ఒకే క్రమం; అంటే ఏ $n$కైనా $\overline{s}(n) = s_{k}(n)$.
ప్రత్యేకంగా $n = k$ తీసుకుంటే, $\overline{s}(k) = s_{k}(k)$ కావాలి.
$s_k(k)$ విలువ $0$ లేదా $1$. అది $0$ అయితే, $\overline{s}(k)$ విలువ
$1$ కావాలి---దాన్ని అలా నిర్వచించాం. $s_k(k) = 1$ అయితే, మళ్లీ
మన $\overline{s}$ నిర్వచనం వల్ల $\overline{s}(k) = 0$.
$\overline{s}$ను ఏ విధంగా నిర్వచించినా రెండు సందర్భాల్లోనూ
$\overline{s}(k) \neq s_{k}(k)$.

$\Bin^\omega$లోని \tetoken{మూలకాలు}{!!{element}s} జాబితా $s_1$, $s_2$, \dots ఉందని
మొదట అనుకున్నాం. ఈ జాబితా నుంచి, అందులో ఉండజాలదని నిరూపించిన క్రమం
$\overline{s}$ను నిర్మించాం. అన్ని $s_i$లు $0$, $1$ల క్రమాలైతే అది
కూడా కచ్చితంగా $0$, $1$ల క్రమమే; అంటే $\overline{s} \in \Bin^\omega$.
ప్రత్యేకించి, $\Bin^\omega$లోని \tetoken{మూలకాలు}{!!{element}s} \emph{అన్నిటి} జాబితా
ఉండజాలదని ఇది చూపుతుంది. అలాంటి ఏ జాబితా నుంచి అయినా, అందులో ఉండదని
హామీ ఉన్న క్రమం $\overline{s}$ను నిర్మించగలం; కాబట్టి అన్ని క్రమాల
$\Bin^\omega$ జాబితా ఉందనే ఊహ విరోధానికి దారితీస్తుంది.
\end{proof}

\begin{explain}
పట్టిక వికర్ణాన్ని వాడి $\overline{s}$ను నిర్వచిస్తుంది కనుక ఈ నిరూపణ
పద్ధతిని ``వికర్ణీకరణ'' అంటారు. వికర్ణీకరణలో తప్పనిసరిగా పట్టిక ఉండాల్సిన
అవసరం లేదు: పట్టిక గానీ, నిజమైన వికర్ణం గానీ లేనప్పుడు కూడా ఇదే
ఆలోచనను వాడి సమితులు \tetoken{లెక్కించదగిన}{!!{enumerable}} కావని చూపవచ్చు.
\end{explain}

\begin{thm}
\ollabel{thm:nonenum-pownat}
$\Pow{\PosInt}$ \tetoken{లెక్కించదగిన}{!!{enumerable}} కాదు.
\end{thm}

\begin{proof}
అదే పద్ధతిలో ముందుకు వెళ్తాం: $\PosInt$ ఉపసమితుల ప్రతి జాబితాకు,
జాబితాలో ఉండజాలని $\PosInt$ ఉపసమితి ఒకటి ఉందని చూపుతాం. కిందిది
$\PosInt$ ఉపసమితుల ఇచ్చిన జాబితా అనుకుందాం:
\[
Z_{1}, Z_{2}, Z_{3}, \dots
\]
ఇప్పుడు ప్రతి $n \in \PosInt$కు $n \in \overline{Z}$ అవడం అంటే
$n \notin Z_{n}$ అయ్యేలా $\overline{Z}$ను నిర్వచిస్తాం:
\[
\overline{Z} = \Setabs{n \in \PosInt}{n \notin Z_n}
\]
ఊహ ప్రకారం ప్రతి $Z_n$, ధన పూర్ణ సంఖ్యల సమితి; కాబట్టి
$\overline{Z}$ కూడా స్పష్టంగా ధన పూర్ణ సంఖ్యల సమితి. అందుచేత
$\overline{Z} \in \Pow{\PosInt}$. కానీ $\overline{Z}$ జాబితాలో
ఉండజాలదు. ఇందుకు ప్రతి $k \in \PosInt$కు $\overline{Z} \neq Z_k$
అని స్థాపిస్తాం.

ఏదైనా $k \in \PosInt$ తీసుకుందాం. ప్రతి $n \in \PosInt$కు
$n \in \overline{Z}$ అవడం అంటే $n \notin Z_n$ అయ్యేలా
$\overline{Z}$ను నిర్వచించాం. ప్రత్యేకంగా $n=k$ తీసుకుంటే,
$k \in \overline{Z}$ అవడం అంటే $k \notin Z_k$. కానీ దీని వల్ల
$\overline{Z} \neq Z_k$: $k$ ఒకదానిలో \tetoken{ఒక మూలకం}{!!a{element}}, మరొకదానిలో కాదు;
అందువల్ల $\overline{Z}$, $Z_k$లకు వేర్వేరు \tetoken{మూలకాలు}{!!{element}s} ఉన్నాయి.
$k$ ఏదైనా కనుక, $\overline{Z}$ అనేది $Z_1$, $Z_2$, \dots జాబితాలో లేదు.
\end{proof}

\begin{explain}
ముందటి నిరూపణ వికర్ణాన్ని ప్రస్తావించలేదు; కానీ దాన్ని ఇలా చిత్రిస్తే
వికర్ణం ఉన్నట్లు భావించవచ్చు. $Z_1$, $Z_2$, \dots సమితులను ఒక
పట్టికలో రాసినట్లు ఊహించండి; ప్రతి \tetoken{మూలకం}{!!{element}} $j \in Z_i$ను $j$వ
నిలువువరుసలో రాయాలి. ఆ జాబితాలోని మొదటి నాలుగు సమితులు
$\{1,2,3,\dots\}$, $\{2, 4, 6, \dots\}$, $\{1,2,5\}$,
$\{3,4,5,\dots\}$ అనుకుందాం. అప్పుడు పట్టిక ఇలా మొదలవుతుంది:
\[
\begin{array}{r@{}rrrrrrr}
  Z_1 = \{ & \mathbf{1}, & 2, & 3, & 4, & 5, & 6, & \dots\}\\
  Z_2 = \{ &  & \mathbf{2}, &  & 4, &  & 6, & \dots\}\\
  Z_3 = \{ & 1, & 2, &  &  & 5\phantom{,} &  & \}\\
  Z_4 = \{ &  &  & 3, & \mathbf{4}, & 5, & 6, & \dots\}\\
  \vdots & & & & & \ddots
\end{array}
\]
అప్పుడు వికర్ణం వెంబడి వెళ్లి, వికర్ణంలో కనిపించే సంఖ్యలను వదిలివేసి,
$j$వ అడ్డువరుస/నిలువువరుసలో ఖాళీ ఉన్న $j$లను చేర్చడం ద్వారా వచ్చిన
సమితి $\overline{Z}$. పై సందర్భంలో $1$, $2$లను వదిలి, $3$ను చేర్చి,
$4$ను వదిలి, ఇదే విధంగా కొనసాగిస్తాం.
\end{explain}

\begin{prob}
వికర్ణ వాదనతో $\Pow{\Nat}$ \tetoken{లెక్కించలేని}{!!{nonenumerable}} అని చూపండి.
\end{prob}

\begin{prob}\label{sfr:siz:nen:prob:f-posint}
ప్రమేయాలు $f \colon \PosInt \to \PosInt$ అన్నిటి సమితి
\tetoken{లెక్కించలేని}{!!{nonenumerable}} అని ప్రత్యక్ష వికర్ణ వాదనతో చూపండి. అంటే,
$f_1$, $f_2$, \dots ప్రమేయాల జాబితా, ప్రతి
$f_i\colon \PosInt \to \PosInt$ ఇచ్చినప్పుడు, ఆ జాబితాలో లేని
ఏదో ఒక $\overline{f}\colon \PosInt \to \PosInt$ ఉందని చూపండి.
\end{prob}

\end{document}
